% =============================================================
%  chapter5_body.tex
%  Shared body (Introduction through Acknowledgments) for Chapter 5,
%  the Draaikraal field study.
%
%  Source and provenance: this is Chapter5/Workspace/Draaikraal 2023 BJT.tex
%  (the raw pandoc export of Draaikraal Report.docx) with informal
%  Word-style headings converted to real \section/\subsection/
%  \subsubsection/\paragraph/\subparagraph commands, each with a
%  \label{sec:ch5-...}, so every part of this chapter is now
%  cross-referenceable via \S\ref{}/\ref{}. No sentence of the original
%  report text was reworded or removed by this conversion -- the only
%  content changes are:
%    - the raw title/author front matter and "Abstract" heading were cut
%      (chapter5.tex's own \chapter{} and \begin{quote} handle those);
%    - the "Works Cited" section was moved out to
%      Chapter5/Workspace/draaikraal_works_cited_raw.tex (not yet
%      integrated with references.bib -- see that file's header);
%    - two bare "Fig N" placeholders and one broken \includegraphics
%      (referencing a media/image1.jpg that was never carried over) were
%      replaced with a real figure (the Draaikraal topographic/geologic
%      map, Chapter5/Figures/source/structural_map.pdf) or \todo{}
%      placeholders, so the chapter compiles cleanly;
%    - two new \section{}s (Introduction, Field Methods) and one new
%      \section{Sample and Station Data} (the xlsx-derived field-data
%      tables) were added as scaffolding, each flagged \todo{} or
%      documented inline where it is editorial rather than original text;
%    - a single new \section{Field Observations} groups the existing
%      Lithologic Descriptions / Structural Observations material
%      (see the NOTE at that heading below).
%
%  Figures for this chapter live in Chapter5/Figures/{source,generated}/
%  (see the \graphicspath entry in dissertation_preamble.tex).
%
%  2026-10-24 update to the newer Draaikraal Report.docx (now in
%  Chapter5/Workspace/): that report differs from the earlier one only in its
%  figures and their captions, which are now all in place (Figs. 1a/1b ->
%  fig:ch5-topo-map, Fig. 2 -> fig:ch5-ma-cleavage-photo, Figs. 3/4 ->
%  fig:ch5-stereonets, Fig. 5 -> fig:ch5-strat-column). The three QGIS maps
%  (lithological, structural, Kliprivier Synform) were added as landscape
%  figures with new captions. The "Sample and Station Data" section moved to
%  the SI (Chapter5/chapter5_SI.tex), and every "photo N" reference now
%  links to the SI photo album via \ChFivePhoto{N}. Pre-update copy:
%  chapter5_body.tex.pre_report_update_20261024.
%
%  2026-10-27: the regional literature review (Layered Mafic Intrusions
%  through The Transformation Hypothesis) moved verbatim to Appendix B
%  (ChapterB/Writing/chapterB_body.tex). What stays here of the old
%  "Geological Background" -- Previous work near Kliprivier and Draaikraal,
%  and Topographic `Inversion' -- is now "Geologic Setting of the Draaikraal
%  Area" and comes before Field Methods. No text changed. Pre-move copy:
%  chapter5_body.tex.pre_litreview_move_20261027; script:
%  Chapter5/Workspace/move_litreview_to_appendix_20261027.py.
% =============================================================

\section{Introduction}
\label{sec:ch5-introduction}

The Bushveld Complex of South Africa in the Mpumalanga, Limpopo, and Gauteng provinces 
hosts the largest layered mafic intrusion in the world (the Rustenburg Layered Series, or RLS), which in 
turn bears the world's largest source of platinum group elements \citep{lee1996} as well as important reserves of chromium, vanadium, titanium, and iron \citep{cawthorn2015}. 
It is exceptionally well-preserved, making it an ideal natural laboratory for studying the geochemistry 
and petrology of layered mafic intrusions. Even so, the order, tempo, and nature of the formation of the
Bushveld Complex remain hotly debated \citep{kruger2020,mungall2016,scoates2021}. Classically, the RLS was thought to have been emplaced as a cohesive, long-lived magma chamber \citep{kruger1994,cawthorn2015},
but recently many have hypothesized that its formation was more complex, with incremental assembly through sill-like intrusions \citep{mitchell2019,yaoetal2021,scoon2023,bourdeau2023}. 
The RLS is divided into the Marginal, Lower, Critical, Main, and Upper zones \citep{vongruenewaldt1985}, 
which are often hypothesized to have formed in distinct events. Thus, the intersections of these stratigraphically
defined zones at the intrusion's margin are particularly informative
locations for detailed field work that seeks to unravel the emplacement history of the RLS. This field study focuses on several
unusual but understudied structural and geomorphic features occurring at
the economically significant intersection of the Main and Critical Zones
with the intrusion's margin in the well-exposed Eastern Limb near the
village of Draaikraal, Mpumalanga, South Africa: the Kliprivier Synform,
the Draaikraal Lineament, and the unexpected recessive topography of the
Main Zone in the agricultural valley surrounding the village of
Draaikraal \citep{vandermerwe2007,molyneux2008}.

\section{Geologic Setting of the Draaikraal Area}
\label{sec:ch5-geological-background}

A broad survey of the literature on the Bushveld Complex is presented in Appendix~\ref{app:bushveld-review}.
Here, the geologic setting in the vicinity of Draaikraal is described.

\subsection{Previous work near Kliprivier and Draaikraal}
\label{sec:ch5-previous-work-near-kliprivier-and-draaikraal}

Kliprivier was the focus of several early works focusing on the margin
of the RLS to better understand the relationship between the RLS and the
sedimentary country rock. \citet{vanbiljon1949} studied the hornfelsized
carbonates of the area and their relationship to the margin, observing
an increase in metamorphic grade to wollastonite skarns near the margin
of the RLS and reporting that the transgressive RLS mafic rocks are
continuous along strike with the metasedimentary rocks, a finding that I
also report here (see \S\ref{sec:ch5-the-transformation-hypothesis}, ``\nameref{sec:ch5-the-transformation-hypothesis}'' for elaboration). He used
this and similar pieces of field and geochemical evidence to develop his
provocative but unpopular Transformation Hypothesis: claiming that the
RLS (and potentially other igneous rocks) were formed from in-situ
magmatic-metasomatic reaction of sediment (see critique and discussion within
the comparatively well-cited work of \citealp{cousins1959}).

Kliprivier and the surrounding De Berg region were studied in detail by
\citet{frick1967}, who distinguished different facies of marginal sills,
documented variable assimilation of country rock by marginal sills and
the presence of Critical Zone and two chromitite seams in the area, and
characterized contact metamorphism within the sedimentary rocks. \citet{frick1967}
emphasized the discordance of the contact between the sediments
and the RLS and reported that chromite seams are disaggregated at the
contact and not present in the Marginal Zone. \citet{button1976} reports that
the extent of the metamorphic thermal aureole decreases southward in the
area as the RLS package thins.

In their influential paper geochemically categorizing the Marginal Zone
sills, \citet{harmer1985} included a broad map of the Marginal
Zone along the Bushveld's Eastern limb in addition to a detailed map of
the marginal rocks at the boundary of the Main and Critical Zones at the
Kliprivier Synform, described more in the \hyperref[sec:ch5-structural-observations]{``Structural Observations''}
section. These maps show that the field area at Draaikraal includes two
distinct classes of Marginal Zone sills. The large map shows ``B2''
gabbroic sills to the north abutting the Critical Zone and ``B3''
gabbroic sills to the south abutting the Main Zone (although later works
argue that the upper Critical Zone continues south within and beyond the
field area: \citealp{vandermerwe2007,latypovetal2022}). Furthermore,
fine gabbros of the ``B2'' and ``B3'' lineages are also subdivided into
``B2N'' and ``B3N''. Because both ``B2N'' and ``B3N'' sills are magnetic
and visually identical in outcrop they are both described here with the
lithologic code ``Rz'', for fine-grained Marginal Zone. Their
coarse-grained ``B2'' gabbros may (in part) correspond to the Disturbed
Gabbronorite (Dcg) lithology mapped in this work.

\citet{vandermerwe2007} mapped and described the geology on Kliprivier and
several nearby farms along the eastern lobe of the RLS to
assess the area's PGE resource potential. He first described the
Kliprivier ``Syncline'', further detailed in the structural
observations, unit descriptions, and discussion in this work. Within, he
references several private studies which were contracted by mining
companies that are not publicly available. Van der Merwe reported a
core log from the area that is helpful to interpret field data.

\citet{latypovetal2022} did a systematic but coarse field study of seven
separate traverses across \textasciitilde60 kilometers of the lower Main
Zone in order to correlate the mottled anorthosites of the lower Main
Zone, arguing that these correlated units, supposedly laterally
continuous across a 4 km change in depth of the RLS floor, imply the
existence of a kilometer-scale melt column. They argue their correlation
requires a voluminous \textasciitilde10\textsuperscript{5}
km\textsuperscript{3} magma chamber during the emplacement of the Main
Zone.

\citet{setera2018} conducted a systematic study of the Main
Zone in the Draaikraal area, reporting major and trace elements of major
cumulus phases. They focused on characterizing magma recharge at the
level of the Pyroxenite Marker in the upper Main Zone, but they also
used chromium concentrations in pyroxene to identify a less significant
magma influx at the level above the mottled anorthosites. This
geochemical cross-section was compared with a similar cross-section to
the north performed in an earlier work by \citet{vantongeren2013},
which notably featured more primitive Main Zone rocks than those
observed near Draaikraal. \citet{setera2018} argue that this
southerly evolving compositional gradient within the Main Zone arose
from evolution of migrating melts away from the feeder of the Bushveld,
interpreted to be at the Thabazimbi-Murchison Lineament.

\subsection{Topographic `Inversion' and the Draaikraal Lineament}
\label{sec:ch5-topographic-inversion-and-the-draaikraal-lineament}

The agricultural village of Draaikraal in the Mpumalanga Province of
South Africa sits in a local sedimentary basin. Geologically, it is
situated within the lower Main Zone of the RLS. The Main Zone is known
to be highly resistant to erosion and dominates surface exposure of the
RLS in the comparatively poorly exposed North and West Limbs of the RLS.
This is thought to be because of the common observation of inversion of
small pigeonite grains into larger, continuous OPX grains, serving as a
sort of resistant ``cement'', contributing to the effectiveness of Main
Zone rocks as dimension stone commonly used as a building material
\citep{cawthorn2015}. Figure~\ref{fig:ch5-topo-map}a shows that this considerable change in
topography at Draaikraal from the northern mountains to the southern
valley occurs despite the continuity of the Main Zone through the
lineament.

% Figs. 1a and 1b of Draaikraal Report.docx (2026-10 version). The 1a caption
% is the report's newer one; the report's older 1a caption (still in the Word
% file just above the figure) contributed the last two sentences of (a).
% Images extracted from the report at full resolution.
\afterpage{%
\begin{landscape}
\begin{figure}[p]
  \centering
  \begin{subfigure}[t]{0.495\linewidth}
    \centering
    \includegraphics[width=\linewidth,height=0.45\textheight,keepaspectratio]{draaikraal_regional_geology_molyneux2008.jpg}
    \caption{}
    \label{fig:ch5-topo-map-geology}
  \end{subfigure}\hfill
  \begin{subfigure}[t]{0.495\linewidth}
    \centering
    \includegraphics[width=\linewidth,height=0.45\textheight,keepaspectratio]{draaikraal_regional_topography_satellite.jpg}
    \caption{}
    \label{fig:ch5-topo-map-satellite}
  \end{subfigure}
  \caption{(a) The field area around Draaikraal with the geological mapping of
    \citet{molyneux2008}. Due North is up on this map. The generally resistant,
    topographically high Steenkampsberg Quartzite and hornfels are mapped as orange
    units on the eastern margin. The Marginal Zone is mapped as the beige color between
    the quartzite and the RLS. The Main Zone is featured in yellow, the Critical Zone in
    the NE quadrant is green and exhibits strikingly recessive topography. A regional
    SW-NE trending fabric is prominent, which agrees with all structural measurements of
    rock cleavage taken in the lower Main Zone and upper Critical Zone aside from the
    rotated units associated with the Kliprivier Synform. The dark box outlines an area
    where SW-NE rock fabrics appear continuous across the Draaikraal Lineament (dashed
    line). The Main Zone around Draaikraal is unusually recessive relative to the Main
    Zone north of the Draaikraal Lineament, or to the west. The compilation of
    \citet{molyneux2008} did not consider more recent mapping that revealed the presence
    of the Upper Critical Zone to the south, terminating beyond the map area
    \citep{vandermerwe2007}.
    (b) The satellite view of the same area as (a), without geology and visual guides
    for clarity. The Draaikraal Lineament remains clearly visible. Elevations are
    colormapped such that relative lows are blue and highs are yellow. The Main Zone is
    generally at higher elevation than the Critical Zone (in the northeast quadrant) and
    the Upper Zone (to the west), while the quartzitic Steenkampsberg Mountains in the
    southeast quadrant are locally the highest feature.}
  \label{fig:ch5-topo-map}
\end{figure}
\end{landscape}}

\section{Methods}
\label{sec:ch5-methods}


\subsection{Field Methods}
\label{sec:ch5-field-methods}
Field observations were made during an 11-day field campaign in July 2023. Field investigation and sampling strategies were
adapted as the study proceeded, but initially were focused on correlating distinct rock units across the Draaikraal Lineament and taking 
structural measurements on each side to characterize it. Geolocated samples were taken across the field area for later geochemical work. 
The observation that plutonic igneous rocks to the north of the Kliprivier Synform were highly unusual and distinct from those to the south
prompted a campaign to more thoroughly characterize its lithologic context and structure. Field observations, samples, and structural measurements
were all geolocated and are recorded in Table~\ref{tab:ch5-field-stations} (Appendix~\ref{app:draaikraal-field-data}). Almost every entry included an interpreted lithological unit to inform
the final geologic map of the area (Figure~\ref{fig:ch5-lithological-map}). Many entries included photographs, which included mostly
close-up photos of outcrops, but also some photos of the surrounding landscape to illustrate the area's geomorphology. Field photographs are cited 
throughout this chapter by number (e.g., photo~\ChFivePhoto{95}); the photographs that
are reproduced in the photo album of Appendix~\ref{app:draaikraal-field-data}
(\S\ref{sec:ch5-SI-photo-album}) are linked from their numbers, and the rest
are in the supplementary photo album. \todo{Link to photo supplement here.}


\subsection{Mapping Criteria}
\label{sec:ch5-mapping-criteria}

The contact between the layered series (Critical and Main Zones) and the
Marginal Zone is often gradational, requiring a somewhat arbitrary
categorization. Since the Lower Main and Upper Critical Zones are
non-magnetic and the finer-grained marginal zone rocks are magnetic, I chose
field-detectable magnetism as the simple discriminant to distinguish
marginal rocks from the massive Main and Critical Zones of the RLS. 

When Layered Gabbronorites (Lg) are
observed in the same outcrops as Mottled or Spotted Anorthosites, the
locality is assigned to be Mottled or Spotted Anorthosite.

\subsection{Structural Measurements}
\label{sec:ch5-structural-methods}

All measurements of plutonic rocks reported here are orientations of
near-planar rock cleavage (\S\ref{sec:ch5-cleavage-layering}).
Strike and dip measurements taken from the Kliprivier Synform were
plotted on a stereonet and analyzed using Stereonet 11 \citep{allmendinger2012,cardozo2013}. Poles to bedding planes of
quartzites, hornfels, mafic sills, and plutonic rocks were calculated
and plotted together in stereonet projection in Figure~\ref{fig:ch5-stereonets}a. In Figure~\ref{fig:ch5-stereonets}b,
the fold axis is calculated using the poles to the bedding planes. The axial plane also was calculated, assuming it
bisected the obtuse angle formed by the limbs. The obtuse fold angle
bisector for the axial plane was selected because the rotation within
the fold observed is quite gentle. An acute bisector for the axial plane
would have been better supported if field evidence showed dramatic changes
in dip along the fold, requiring a complete structural overturn that
locally annihilated the original horizontality of the deformed rocks in
addition to steeply dipping rocks intermediate to the two limbs.

\subsection{Stratigraphic Column Analytical Methods}\label{sec:ch5-stratigraphic-column-methods}

The field data of this study north of the Kliprivier Synform are
projected onto a stratigraphic column following the assumption that all
rocks in the field area north of the Kliprivier Synform follow
Steno's Law of Original Horizontality and were only uniformly deformed
by tilting due to isostatic flexure towards the center of the RLS after
emplacement, so that all regional bedding and structures are currently
dipping 15 degrees to the northwest relative to their original
orientations. Geolocated field data are projected onto a line to the north of the
Kliprivier Synform designed to be orthogonal to the regional strike
(trending 301.79 degrees: W-NW) using the `Snap Geometries to Layer' function in
the `Vector Geometry' Processing Toolbox in QGIS, with a tolerance of 1
degree set to ``prefer closest point''. Section heights then were calculated for
these projected points along the reference line to an arbitrary point along this line,
to get relative stratigraphic heights using the  ``Distance Matrix'' function 
in the ``Vector Analysis'' Processing Toolbox. A correction was added to these lateral distances to
account for elevation, calculated as the product of the elevation
differential and the tangent of the regional dip, assumed to be 15°. The
column heights are normalized to the height of the Merensky Reef. Because
the Merensky Reef was not observed in the field area, this was instead done by proxy: stratigraphic heights were
calculated relative to the UG2 chromitite observed in the field area (\S\ref{sec:ch5-unit-ug2}) at a fixed height of -86 meters. 
This chromitite is mapped by \citet{vandermerwe2007} as the UG2, who profiled drill cores 
and determined the Merensky Reef to be 86 meters above the UG2 in the field area \citep{vandermerwe2007}. Note that, for simplicity, 
it is assumed that there are no local changes in strike and dip,
even though the dip of the Main Zone likely shallows and rotates to be
roughly due east further from the margin of the intrusion. This approximation is likely to be acceptable locally for the upper 
Critical and lower Main zones, but becomes less accurate further from the margin. Therefore, the
stratigraphic heights presented here for the Homogeneous Gabbronorite
lithology far upsection are likely overestimated. The interpreted schematic
stratigraphic column for the region is graphically presented in
Figure~\ref{fig:ch5-strat-column}.

\begin{figure}[htbp]
  \centering
  \includegraphics[width=\linewidth]{draaikraal_schematic_strat_section.jpg}
  \caption{The schematic stratigraphic cross section of the RLS in Draaikraal. The
    Kliprivier Synform is represented by the slab of floor rocks protruding into (and
    creating) the Disturbed Critical Zone at the halfway point of the column. The
    Draaikraal Lineament would form from upward percolating fluids trapped and channeled
    by this protruding slab of floor rocks in a line directly above the slab in this
    view. The Kliprivier Synform is not interpreted to have deformed the RLS above the
    UG2.}
  \label{fig:ch5-strat-column}
\end{figure}

\section{Results: Field Observations}
\label{sec:ch5-field-observations}

Observations are organized into three sections: lithologic descriptions, rock cleavage and modal layering, and structural observations; the first and last each have their own map.

\subsection{Lithologic Descriptions}
\label{sec:ch5-lithologic-descriptions}

Lithologic descriptions are organized by the interpreted map units (see Figure~\ref{fig:ch5-lithological-map}) and their constituent lithologies. 
Three mapping units are described in turn, each with a wide variety of lithologies: the marginal rocks (the floor and walls of the intrusion, 
including sedimentary country rock and concordant pre- and syn-Bushveld magmatic sills); the plutonic rocks of the RLS;
and secondary, post-Bushveld rocks that were largely ancillary to the field study. 

\afterpage{%
\begin{landscape}
\begin{figure}[p]
  \centering
  \includegraphics[width=\linewidth,height=0.8\textheight,keepaspectratio]{lithological_map.pdf}
  \caption{Lithological map of the Draaikraal field area, eastern limb of the Bushveld
    Complex. Field observations of outcrop are colored by lithology over the mappable units (floor,
    Disturbed Critical Zone, Layered Critical Zone, Transitional Unit, Main Zone). The
    lower Main Zone anorthosite marker beds (LMA, MMA, UMA, and the spotted
    anorthosites) are traced in color, and the Draaikraal Lineament is dashed. Field observation
    data are listed in Table~\ref{tab:ch5-field-stations}.}
  \label{fig:ch5-lithological-map}
\end{figure}
\end{landscape}}

\subsubsection{Margin}
\label{sec:ch5-eastern-margin}

Much of this study focused on the eastern margin of the eastern limb of
the Rustenburg Layered Series (RLS) east of Draaikraal. The margin
displays the most unambiguous bedding planes of the study area,
generally dipping 25° to the northwest towards the center of the
Bushveld Complex, except for the Kliprivier Synform, which is detailed
below in the \hyperref[sec:ch5-structural-observations]{``Structural Observations''} section. The marginal rocks
within the field area are composed of gabbroic intrusive igneous sills
concordantly intercalated with variably metamorphosed, melted, or
otherwise devolatilized Transvaal sediments of the Steenkampsberg and
Houtenbek formations. Occasionally, melting and/or reactive porous flow
of melts through sedimentary layers remineralizes sedimentary layers
with igneous mineralogy and texture (photos~\ChFivePhoto{35}, \ChFivePhoto{51}, and \ChFivePhoto{52}). Concordant
sills of marginal zone magmas and packages of sediment range in
thickness from 10 to 100s of meters, and can be found mutually
interbedded at least a kilometer from the margin of the RLS. Here, the
marginal rocks are divided into three units: ``Rz'', the fine-grained
Marginal Zone gabbroic sills/magmas; ``Qz'', relatively undisturbed
Houtenbek quartzite; and ``ScP'', hornfelsed metapelites, marbles, and
quartzites with related igneous/hydrothermal pegmatites.

\subsubsection{Rz -- Fine-grained gabbro}
\label{sec:ch5-unit-rz}

The Marginal Zone gabbroic sills are commonly fine grained (\textless1
mm), but locally are coarser with some grains up to 2 mm in the
groundmass. Unlike much of the plutonic RLS, the Marginal Zone is not
visibly modally layered. Marginal Zone magmas within a meter of
Steenkampsberg Quartzite often contain conspicuous evidence of silicic
contamination: although they retain an igneous texture, the color
significantly lightens and the rocks become much harder and less
magnetic approaching the contacts with Steenkampsberg Quartzite. Aside
from these immediate contacts, the Marginal Zone gabbro is always
magnetic. The Marginal Zone contains more pyroxene (distinguishing
clinopyroxene and orthopyroxene is difficult with such small crystals)
than plagioclase with variable magnetite, sometimes up to 10\%, and
minor phlogopite. The heavy, fine fraction of mineral separates from one
sample was briefly studied under optical microscope to identify trace
phases, which include apatite, zircon, sulfides, and euhedral chromite.

The Marginal Zone magmas contain igneous xenoliths and are rich in
paleo-vertical, across-strike features that likely record rising,
buoyant fluid/magma. Perfect planar resistance banding (photo~\ChFivePhoto{95}) is the
most common structure in the fine grained Marginal Zone and is nearly
ubiquitous within the unit. Though the resistance bands have seemingly
random strikes, they most often are steeply dipping and may represent
fluid conduits along brittle fractures within the Marginal Zone.

In previous works, isolated meta-sedimentary xenoliths \citep{buick2000}
were emphasized, but in this study, meta-sedimentary rocks were
never found as minor xenoliths within magnetic marginal sills, but
rather as continuous, concordant strata (photos~\ChFivePhoto{42}, \ChFivePhoto{70}, \ChFivePhoto{71}) or as
locally abundant disaggregated broken strata cemented within a pink,
evolved diatexite (photo~\ChFivePhoto{100}), both designated as a distinct unit: ScP.

The magnetic sills of the Marginal Zone (Rz) did include enclaves of
coarse-grained leucogabbronorite, mottled anorthosite or pyroxenite
(photos~\ChFivePhoto{2},\ChFivePhoto{4},\ChFivePhoto{6},\ChFivePhoto{96}). The anorthosite-rich enclaves are rounded, 6-10 cm
along their longest dimension, and are imbricated to be preferentially
lying flat paleohorizontally within the Marginal Zone, suggesting they
were weathered/eroded by reaction and/or mechanical transport through
Marginal Zone magma. Pyroxenitic enclaves are also often spherical
(photos~\ChFivePhoto{4}, \ChFivePhoto{96}) and may be related to the formation of the pyroxenite
``Tennis Ball Marker'' (TBM) of the Main and Critical Zones, described
further in the discussion. These plagioclase-rich igneous xenoliths
contain large (\textless3 cm) optically continuous oikocrystic augite
that enclose blocky plagioclase (\textless5 mm).

More commonly observed within the magnetic Marginal Zone sills east of
Draaikraal are paleovertical structures that suggest buoyant upwelling
of fluids and melts sourced from the hornfelsized sedimentary rocks.
These structures included sheared, viscously deformed lenses of
thin pegmatites, including anorthosite (photo~\ChFivePhoto{6}), magnetite gabbros
(photo~\ChFivePhoto{60}, sheared gabbro \textasciitilde{}15 cm;
Figure~\ref{fig:ch5-photo60-deformation}), and evolved
quartz-biotite pegmatites (photo~\ChFivePhoto{67}: 1$\times$3~cm). In addition to homogeneous
pegmatites, the Marginal Zone gabbros host paleovertical invasions of
anastomosing anorthosite derived from anatexis of underlying
metasedimentary rocks (photos~\ChFivePhoto{6}, \ChFivePhoto{97}). These leucosome intrusions
occasionally carry blocks of meta-sedimentary paleosome or brecciate and
enclose the Marginal Zone gabbronoritic host rock itself (photo~\ChFivePhoto{98}).

Near the contact with the migmatite, differential weathering sometimes
reveals proud, thin, pinched ovate paleovertical structures (5$\times$1~cm
scale) on weathered faces of Marginal Zone norite that are otherwise
invisible on a fresh rock face {[}photo~\ChFivePhoto{40}{]}.


Fine-grained Marginal Zone gabbros are mapped up to 120 meters below the
Merensky Reef stratigraphically.

\subsubsection{Qz -- Houtenbek Quartzite}
\label{sec:ch5-unit-qz}

The Houtenbek quartzite dips concordantly with marginal zone sills and
migmatized sedimentary rocks towards the center of the Bushveld, here
dipping about 25° to the northwest. Even within meters of a contact with
Marginal Zone sills (Rz), the quartzite often retains a clastic
sedimentary texture and possibly could be considered a proper
clast-dominant quartzose sandstone with well-sorted and well-rounded
coarse quartz grains \textasciitilde1 mm in diameter with no other
mineral visible. While there is little structure visible internally
because of the purity of the quartzite, the bedding planes occasionally
reveal conspicuous asymmetric ripple structures {[}photos~\ChFivePhoto{2}, \ChFivePhoto{68}{]}
\textless1 cm in height. One in-place ripple structure suggests a
paleoflow direction towards the south-southwest. Some quartzites are
hornfelsized: recrystallized into an interlocking crystalline texture
with large, irregular, angular, coarse grains \textless2 cm across. Many
quartzites with preserved clastic textures near the margin of the RLS
are stained bright red within rocks' unweathered interior. This staining
is hypothesized to be derived from mobilized Fe-rich fluids percolating
along grain boundaries, probably caused by heat and/or fluid circulation
related to the emplacement of the RLS. Steenkampsberg Quartzite is mapped
up to 140 meters below the Merensky Reef stratigraphically.

\subsubsection{ScP -- Houtenbek Hornfels, Migmatites, Pegmatites}
\label{sec:ch5-unit-scp}

In addition to quartzite, the Houtenbek formation contains carbonate and
pelitic horizons \citep{button1976}. At the margin east of Draaikraal, there
is a 60 m thick package of melted and metamorphosed sediment containing
stratiform pegmatites with variable mineralogy. Hornfelsized sediments
come in two main varieties near Draaikraal: pink, diatexite-supported
breccias {[}photo~\ChFivePhoto{100}{]} or dark, high-grade schists {[}photo~\ChFivePhoto{29}{]} with
thin but conspicuous navy, gray, white, tan, and pink bands of
relatively mafic and felsic minerals that have been recrystallized and no
longer bear clastic textures except for thin seams of clast-dominated
quartzite in a low-fraction glassy biotite-rich matrix {[}photo~\ChFivePhoto{35},
large light seam is \textasciitilde4 cm thick{]}. This unit is rarely
magnetic and most often is visible as continuous, concordant strata
{[}photo~\ChFivePhoto{42}{]}. Even when highly brecciated, this unit can be followed
along the strike of the quartzites and Marginal Zone, implying that the
sediments were not incorporated into Marginal Magmas, but that they
rather were melted and devolatilized in-situ due to heating from the
emplacement of adjacent marginal sills. Fluids and melts produced in
this process are visible as a wide variety of stratiform pegmatites and
skarns with undulating erosive contacts within comparatively planar
meta-sediments {[}photo~\ChFivePhoto{30}: feldspar/wollastonite(?) + amphibole skarn;
photo~\ChFivePhoto{69}: magnesite, quartz pegmatite with minor siderite{]} and as
pockets following fluid flow pipes and void spaces that are accommodated
by the host strata through ductile deformation {[}photo~\ChFivePhoto{37}, \ChFivePhoto{38}:
Orthopyroxene and quartz pegmatite{]}. Hornfels in the field area are
mapped up to 190 meters below the Merensky Reef stratigraphically.

\subsubsection{The Rustenburg Layered Series (RLS)}
\label{sec:ch5-rls-units}

The Rustenburg Layered Series is the world's largest individual igneous
intrusion \citep{cawthorn2015,latypovetal2022} and is subdivided into the Lower Zone, Critical Zone, Main Zone, and Upper Zone (\S\ref{sec:ch5-the-bushveld-event-and-rls}). In this study, RLS
rocks of the lower Main Zone, Transitional Zone, and upper Critical Zone
are mapped. The rocks of the Marginal Zone (Rz) are not layered, nor
considered here to be a part of the RLS on a textural basis. Marginal
zone rocks weather distinctly from the plutonic RLS phase with abundant,
sharp, \textasciitilde3 cm-scale potholes and generally are more
textured and rough {[}photo~\ChFivePhoto{36}{]} on a weathered surface than RLS rocks,
which weather spheroidally and proudly {[}photos~\ChFivePhoto{8}, \ChFivePhoto{11}: ``porphyritic''
gabbronorite{]}. Xenolith blocks of Marginal Zone norites bearing their
distinctive resistance banding are found in both the Critical and the
Main Zones {[}Critical Zone: photo~\ChFivePhoto{76}; Main Zone: \ChFivePhoto{91}{]}, though each
xenolith sighting within the Main Zone is proximal to a mottled or
spotted anorthosite layer. These xenoliths are variably annealed,
recrystallized, and even ductily deformed {[}Critical Zone: photo~\ChFivePhoto{77};
Main Zone: \ChFivePhoto{84}{]}.

The RLS is famous for its distinctive and numerous laterally continuous
stratiform marker horizons, sometimes stretching on a 100-km scale
\citep{cawthorn2015}, which include the Critical Zone chromitite cycles,
Merensky and Bastard pegmatoid reefs, three mottled anorthosite packages
in the lower Main Zone, the Pyroxenite Marker of the upper Main Zone,
and the magnetitite cycles of the Upper Zone. In this study, I observe
and describe the three lower Main Zone magnetic mottled anorthosites
{[}photos~\ChFivePhoto{49}, \ChFivePhoto{53}, \ChFivePhoto{86}, \ChFivePhoto{93}, \ChFivePhoto{105}, \ChFivePhoto{106}, \ChFivePhoto{108}, \ChFivePhoto{111}{]}, the magnetic chromitite
UG2 {[}photos~\ChFivePhoto{65}, \ChFivePhoto{66}{]}, and a non-magnetic upper Critical Zone mottled
anorthosite that is probably a part of the Merensky Reef footwall
{[}photo~\ChFivePhoto{81}{]}, and several packages of xenolith-bearing modally layered
gabbronorites and spotted anorthosites within the Main Zone (photos~\ChFivePhoto{91} and \ChFivePhoto{111}).
These laterally continuous reefs are mapped and described separately
from the bulk cumulate rocks of the RLS.

\subsubsection{Cz: Critical Zone}
\label{sec:ch5-zone-cz}

The Critical Zone was previously mapped to pinch out due east of
Draaikraal \citep[shown in the compilation of][]{molyneux2008}; however,
many now recognize that a thin layer of the upper Critical Zone persists
much further south towards Tontledoos \citep{vandermerwe2007}. In this
study, diagnostic Critical Zone lithologies were only observed north of
the Kliprivier Synform, where typical finely modally layered norites
and anorthosites are well-exposed (photos~\ChFivePhoto{64} and \ChFivePhoto{80}), as well as the UG2
chromitite. The Critical Zone above the Kliprivier Synform dips
\textasciitilde15° to the northwest similarly to the majority of the
Main and Marginal Zone rocks in the field area, but gradually rotates to
a southwestern dip as it approaches the contact with the rotated block
of metasediments and gabbronorite sills of the Kliprivier Synform. It is
essential to note that some rotated Critical Zone rocks are structurally
beneath the large, rotated package of intruded metasediments (Rz + Qz +
ScP) in the area of the Kliprivier Synform, meaning that this package of
country rock is not only serving as the floor for RLS rocks to the west,
but is also (at least locally) a roof for Critical Zone rocks below the
level of the UG2.

\subparagraph{Dcg: Disturbed Gabbronorite}
\label{sec:ch5-unit-dcg}

This lithology is restricted to the immediate area of the Kliprivier
Synform. Disturbed Gabbronorites are coarse-grained, even pegmatoidal
(photo~\ChFivePhoto{62}), with feldspars and pyroxenes up to 3~cm long. This unit was
named in response to \citet{vandermerwe2007}, who said that the
``structural change {[}of the Kliprivier Synform{]} is reflected in the
Bushveld rocks, which become highly disturbed and fragmented''.
Disturbed Gabbronorites are most recognizable for their abundance of
blobby autoliths of anorthosite (photo~\ChFivePhoto{61}) and pyroxenite (photos~\ChFivePhoto{70} and
\ChFivePhoto{72}). While it exhibits abundant examples of modal separation, they are
never remotely planar, unlike the majority of the RLS. 

This lithology is most often non-magnetic, but a few specimens were
magnetic. Besides the mottled anorthosites of the Main Zone and the
anorthosite-hosted chromitite leader of the UG2, no other rocks of the
RLS within the field area are magnetic. Many but not all Disturbed
Gabbronorites are rotated within the Kliprivier Synform. Generally, the
Disturbed Gabbronorites are CPX-poor and sulfide-rich with abundant,
euhedral \textless1 cm pyroxenes set in an anorthosite matrix that
effectively isolates the large, euhedral pyroxenes from one another
{[}photo~\ChFivePhoto{62}{]}. This unit generally becomes more feldspathic as one
moves away from the intruded metasediments of the Kliprivier Synform
into the Critical Zone. This lithologic unit contains recrystallized and
annealed xenolith blocks {[}photos~\ChFivePhoto{74}, \ChFivePhoto{75}, and \ChFivePhoto{77}{]} of fine-grained
Marginal Zone gabbros (Rz) as well as coarse-grained pockets containing
very large (\textless{} 5cm) oikocrystic pyroxene {[}photo~\ChFivePhoto{73}{]} and
large clusters of pyroxene often rimmed by thick anorthosite halos
{[}photo~\ChFivePhoto{72}{]}. Disturbed Gabbronorite is mapped between 120-145 meters
below the Merensky Reef stratigraphically, around the level of the UG1
in core logs \citep{vandermerwe2007}. \citet{vandermerwe2007} and
\citet{latypovetal2022} report isolated patches of massive chromite here, but I did
not see any.

\subparagraph{UG2: Upper Group 2 Chromitite}
\label{sec:ch5-unit-ug2}

Near the top of the upper Critical Zone in the extreme northern portion
of Kliprivier farm, there is a small but high-quality exposure of
chromitite (\textgreater30 cm thick). The bulk of the exposure is a
melanocratic, non-magnetic, anorthositic chromite, which is modally
\textasciitilde50\% \textless1mm euhedral chromite set within a fine
matrix of \textless1 mm plagioclase (photo~\ChFivePhoto{66}). Structurally above this
massive layer are some thin (\textless1 cm thick) chromitite leaders
hosted within bright white mottled anorthosite (photo~\ChFivePhoto{65}). Both layers
contain amoeboid oikocrystic pyroxene-olivine mottles (\textless2cm).
The mafic silicate mottles and anorthosite matrix seem to have
crystallized within the interstices of a preexisting framework of a
continuous chromite crystal slurry {[}photos~\ChFivePhoto{65}, \ChFivePhoto{66}{]}, \citep[see][for similar interpretations]{voordouw2009,mukherjee2017}. The
massive melanocratic layer includes pure (relatively chromite-free)
white anorthosite blocks (\textless4$\times$15~cm) that are lying flat and
dipping west towards the center of the Bushveld Complex. These pure
anorthosite blocks contain thin chromitite stringers that together have
been ductilely deformed, and a thin (\textasciitilde1 mm) white seam of
pure, apparently undeformed anorthosite forms a plane bounding the top
of some of these blocks, also dipping to the west.

The anastomosing chromite leaders within the overlying mottled
anorthosite are sometimes highly magnetic and are a much deeper,
vitreous black: appearing under hand lens to be \textgreater80\% modally
chromite/magnetite {[}photo~\ChFivePhoto{65}{]}.

This chromitite package is interpreted as UG2 for a few reasons. First,
although the anastomosing chromite leader resembles the UG1 stringers,
the leader is structurally above the massive chromitite, whereas the UG1
stringers are recognized below the massive chromitite \citep{voordouw2009,mukherjee2017}. The UG2 is known to have several small
``leaders'' a small distance above the main seam \citep{maier2024}. Second, the magnetic chromitite within the leader bears
explaining. One heterogeneous magnetic chromitite at the top of the UG2
with discrete magnetic and non-magnetic bands was described by
\citet{frankel1942}, 20-30 km to the north of Kliprivier at Tweefontein.
Third, a rock at this location was also mapped as the UG2 by \citet[Fig.~4]{harmer1985} and \citet[Fig.~3]{vandermerwe2007}, who also
identified the UG2 and closely overlying mottled anorthosites in drill
cores of the area using PGE grades. \citet{vandermerwe2007} also reports
that the UG2 at Kliprivier is unusually enriched in iron, which may
explain the magnetism of the leader. The UG2 is 86 meters below the
Merensky Reef in the direct vicinity of the Kliprivier Synform \citep[Borehole~3:][]{vandermerwe2007}.

\subparagraph{Ln: Layered Norite:}
\label{sec:ch5-unit-ln}

Spectacular outcrops of modally layered norites within the upper
Critical Zone are observed within the field area at Kliprivier around
and above the stratigraphic level of the UG2 chromitite (photos~\ChFivePhoto{64} and
\ChFivePhoto{80}), at stratigraphic heights from 100 to 50 meters below the Merensky
Reef. Modal layering is pronounced, with pyroxenite-rich layers being
separated by near-pure centimetric anorthosite bands. This layering is
undulose and irregular.

\subparagraph{Px: Pyroxenite}
\label{sec:ch5-unit-px}

This refers to rocks predominantly composed of pyroxenes. Pyroxenites
are rare in the field area. There was one coarse-grained sulfide-bearing
feldspathic pyroxenite horizon in the Critical Zone just below the UG2
(photo~\ChFivePhoto{78}) as well as a pyroxenite layer just below a mottled
anorthosite layer in the UMA in the lower Main Zone (photos~\ChFivePhoto{53}--\ChFivePhoto{54}).

\subsubsection{Tz: Transitional Zone}
\label{sec:ch5-zone-tz}

There is no perfect division for the Main and Critical Zones. Moving
upsection from the Critical Zone into the Main Zone, Cr\textsubscript{2}O\textsubscript{3} concentration within OPX decreases, strontium isotopes become more
radiogenic, and cumulus clinopyroxene becomes important--- but these
changes occur at different stratigraphic heights within the transition
zone \citep{cawthorn2015} between the Merensky Reef and Giant Mottled
Anorthosite (or above). Within the literature, the bottom of the Main
Zone is considered to be above the Giant Mottled Anorthosite at the top
of the Bastard unit, and the Merensky reef is taken to be the top of the
Upper Critical Zone \citep{cawthorn2015,mitchell2002,vandermerwe2007}. The Bastard and Merensky pyroxenites and Giant Mottled
Anorthosite were never observed in the field area, but a
\textasciitilde50 cm thick mottled anorthosite {[}photo~\ChFivePhoto{81}{]} was
observed \textasciitilde50 meters upsection from the UG2 with only
finely layered norite {[}photo~\ChFivePhoto{80}{]} and porphyritic gabbronorite
outcropping between the UG2 and the mottled anorthosite. Since
\citet{vandermerwe2007} reports in core logs that the Merensky Reef is 86 meters
above the UG2 on Kliprivier and \citet{latypovetal2022} observe unusual
pyroxenites just above this mottled anorthosite, this anorthosite layer
is interpreted to be a part of the footwall to the Merensky Reef, and
the pyroxenites described in \citet{latypovetal2022} are interpreted to be
the Merensky Pyroxenite.

\subparagraph{Pg: Porphyritic Gabbronorite}
\label{sec:ch5-unit-pg}

The ``Porphyritic Norite'' is a well-recognized marker of the lower Main
Zone, with large, rounded pyroxene phenocrysts set in a white,
feldspar-rich matrix, used often in mining to locate the Merensky Reef
\citep{sacs1980,vongruenewaldt1985}. I observe a general textural
similarity in lower Main Zone Gabbronorites with the unlayered upper
Critical Zone gabbronorites below the Merensky Reef (photos~\ChFivePhoto{10}, \ChFivePhoto{79}).
These gabbronorites are distinct from the dark and homogeneous
gabbronorites that compose the majority of the Main Zone (Hg). The
porphyritic gabbronorite contains large (\textasciitilde\textless{} 1cm),
equant pyroxenes set in a leuconoritic matrix that is predominantly
plagioclase with small spots of interstitial orthopyroxene. I refer here
to them collectively as the ``porphyritic gabbronorite'' because I see
cumulus clinopyroxene present throughout, becoming more modally abundant
upsection until about half of the cumulus pyroxenes are CPX near the
Lower Mottled Anorthosite (LMA). Generally, cumulus CPX is present in
the Main Zone and absent from the Critical Zone \citep{cawthorn2015}.
However, some cumulus CPX has been reported at the top of the Critical
Zone \citep{cameron1982,sacs1980}.

The rocks of this unit have well developed cleavage {[}photo~\ChFivePhoto{63}{]};
however, they are not visibly layered. I located porphyritic
gabbronorite south of the Kliprivier Synform at the beginning of the
second field day and assumed it was Main Zone given the cumulus CPX I
saw, but now I believe that the porphyritic gabbronorite occurs in the
upper Critical Zone through the Transition Zone and into the lower Main
Zone in the vicinity of Draaikraal. The Porphyritic Gabbronorite in this
field area hosts spherical pyroxenite tennisballs (photos~\ChFivePhoto{10}, \ChFivePhoto{79}), in
the same locations as mapped by \citet{harmer1985} and
\citet{mitchell2023}. Porphyritic Gabbronorite is mapped in the Critical and
Main Zones 480 meters below the Merensky Reef to 380 meters above,
stratigraphically.

\subsubsection{Mz: Main Zone}
\label{sec:ch5-zone-mz}

The Main Zone is the most volumetrically significant portion of the RLS:
as thick as 5 km \citep{latypovetal2022}. The Main Zone stratigraphy is
(in general) relatively featureless aside from three prominent mottled
anorthosite packages in the lower Main Zone and the Pyroxenite Marker at
the upper Main Zone, which was not studied in this work. Mottled
anorthosites and associated strongly modally layered sequences are
mapped and described here in order to study the effect of the Kliprivier
Synform and Draaikraal Lineament on the stratigraphy, as these unique
layers can be traced great distances through the Lower Main Zone
\citep{latypovetal2022}.

\subparagraph{Hg: Homogeneous Gabbronorite:}
\label{sec:ch5-unit-hg}

The massive Homogeneous Gabbronorite of the Main Zone varies slightly in
color due to variations in the modal abundance of pyroxenes and
plagioclase, though these variations are subtle and continuous on a
scale of 10s of meters. There is no obvious fabric visible in outcrops
or hand samples of these massive rocks, although I was shown a polished
core of upper Main Zone gabbronorite from a dimension stone mine east of
Roossenekal that obviously displayed a fabric where crystals
preferentially were aligned with the gravitational plane. Stonemasons
favor this rock for building, but are careful to allow only faces cut
against this fabric to outwardly show for aesthetics and durability.
Homogeneous Gabbronorite (Hg) begins about 450 meters above the Merensky
Reef and continues upsection stratigraphically through the Pyroxenite
Marker \citep{sacs1980,vongruenewaldt1985}.

\subparagraph{(L/M/U)MA: (Lower/Middle/Upper) Mottled Anorthosites}
\label{sec:ch5-unit-lmuma}

Mottled Anorthosites are enigmatic igneous rocks composed of a nearly
pure, white anorthosite matrix containing (often quite large: \textless5
cm) mossy green ``mottles'' of oikocrystic pyroxene +/- olivine that
enclose abundant plagioclase. In the RLS they are found in large
stratiform horizons interpreted to be continuous across the majority of
the Main Zone in the Bushveld (\textgreater100 kilometers:
\citealp{latypovetal2022}). The lower Main Zone near Draaikraal contains 3 large mottled
anorthosite packages. These mottled anorthosites are magnetic and
notably contain rare large (\textless{} 2mm) grains of magnetite within
the mottles. They are found in association with larger packages of
distinct but non-magnetic anomalous lithologies. The associated unusual
rocks include: 1) fine green gabbronorites with blebby, anhedral
phenocrystic (\textless2 cm) clinopyroxene (lithologic code Fgog:
Fine-grained oikocrystic gabbronorite) associated with the UMA; 2)
feature-rich, coarse pyroxenites (lithologic code: Px, {[}photos~\ChFivePhoto{53},
\ChFivePhoto{54}{]}) and 3) blobby pyroxenitic gabbronorites (lithologic code Bap:
Blobby anorthosite-pyroxenite) that include deformed, irregular,
decimeter-scale blobby pyroxene agglomerations within an anorthositic
matrix that are interpreted here to be recrystallized blocks {[}photos~\ChFivePhoto{104}, \ChFivePhoto{108}{]}. Sometimes these blobby anorthosite-pyroxenites exhibit
modal layering of coarse (\textless2 cm) pyroxene oikocrysts {[}photo~\ChFivePhoto{105}{]} within anorthosite. The mottled anorthosites and their associated
pyroxenites and blobby gabbronorites have undulating, nonplanar contacts
with one another {[}photo~\ChFivePhoto{49}{]} and with the underlying common
gabbronorites {[}photos~\ChFivePhoto{107}, \ChFivePhoto{103}{]}. Sometimes these contacts may be
drastically discordant from the local strike and dip. In photo~\ChFivePhoto{106}, the
mottled anorthosite - gabbronorite contact is dipping \textasciitilde70°
to the east towards the margin of the chamber: a much steeper dip and
opposite strike than observed in the leucogabbronorites of the immediate
vicinity. Occasional 3D exposure of the anorthosite-gabbronorite contact
allowed me to measure the contact rotating 90+ degrees in this outcrop
within meters (photo~\ChFivePhoto{107}). It is unclear if these undulating contacts
are evidence for discordance of the mottled anorthosites, since in other
places the contact appears concordant {[}photo~\ChFivePhoto{110}{]}, and local
apparent discordance could be the result of thermal erosion of the
cumulate pile by overlying magmas \citep{roberts2007}.

 

In this study, all three lower Main Zone anorthosites described by
\citet{latypovetal2022} and mapped by \citet{vandermerwe2007} were located
within the low-lying valley south of the Draaikraal Lineament at
stratigraphic heights between 200 and 600 meters above the level of the
Merensky Reef. The Upper and Middle ``mottled'' anorthosites in the
mountains to the north reported by \citet{latypovetal2022} were located in
this study, but are visually quite distinct and occupy higher
stratigraphic heights, argued here to be
distinct units that stratigraphically overlie the thick, massive mottled
anorthosites in the valley to the south (see the Discussion of this
work). \citet{vongruenewaldt1985} report multiple subordinate
anorthositic layers near and at the boundary of the Main Zone subzone A
and B much higher above the level of the Main Mottled Anorthosite. Those
higher anorthosites are referred to here as spotted anorthosites (SA) to
reflect their distinct texture. 

\subparagraph{SA: Spotted Anorthosites}
\label{sec:ch5-unit-sa}

North of the Draaikraal Lineament where exposure is better, two packages
of anorthosites were found in the lower Main Zone. These anorthosites
are somewhat distinct from the common, unlayered gabbronorites of the
Main Zone, but anorthositic beds are significantly smaller (\textless1 m
thick) and altogether less visually striking than the Mottled
Anorthosites (MA) unit described above. The spotted anorthosites are
non-magnetic and bear \textless1 cm spots of equant pyroxenes composing
\textless20\% of the rock by volume {[}photos~\ChFivePhoto{86}, \ChFivePhoto{87}{]} that sometimes
contain recrystallized pyroxenite-gabbronorite blocks {[}photo~\ChFivePhoto{111}{]}.
The Spotted Anorthosites are strongly associated with much thicker
(\textless10 m) packages of planar, finely layered gabbronorites that
are unusual within the Main Zone and are more reminiscent of the
Critical Zone {[}photos~\ChFivePhoto{85}, \ChFivePhoto{112}{]}. These modally layered sections also
include xenoliths {[}photos~\ChFivePhoto{91}, \ChFivePhoto{84}{]}, which were only observed within
the Main Zone in these modally layered packages. These Spotted
Anorthosite packages are argued here to be distinct and
stratigraphically above the remarkable and prominent Mottled
Anorthosites described above. Two or three packages of Spotted
Anorthosites (often including thick layers of Layered Gabbronorite) are
interpreted here at stratigraphic heights 700-1050 meters above the
Merensky Reef. More packages may exist, although they are unlikely to be
exposed to the south of the Draaikraal Lineament.

\subparagraph{Lg: Layered Gabbronorite}
\label{sec:ch5-unit-lg}

Modal layering is generally less common in the Main Zone than the
underlying Critical Zone, and in outcrop it is comparatively
understated. Whereas the Ln (Layered Norite) unit of the Critical Zone
contains vivid, centimeter-scale, near-monomineralic anorthositic bands,
the modal layering in the Main Zone gabbronorites is millimeter-scale,
represented by subtle differences in cumulus grain size and modal
proportions, and the contacts between layers are planar (photos~\ChFivePhoto{85} and
\ChFivePhoto{112}). The Layered Gabbronorites are associated with Mottled Anorthosites
(MA) or Spotted Anorthosites (SA).  Layered
Gabbronorites are mapped here at stratigraphic heights between 400 and 1050 meters in association with Mottled and Spotted Anorthosites.



\subsubsection{Post-Bushveld Rocks}
\label{sec:ch5-post-bushveld-rocks}

In the valley surrounding Draaikraal, alluvium primarily is composed of
well-rounded quartzite cobbles (\textasciitilde10 cm) and ironstone
laterites (photo~\ChFivePhoto{59}). The ironstone laterites are commonly found as both float in
alluvium and as in-place outcrops. Bushveld mafic rocks are rare in
alluvium, probably due to their relative susceptibility to chemical
weathering on Earth's surface. Quartzite was substantially dominant over
RLS mafic rocks in alluvium even 8 km away from the nearest known
quartzite source within the RLS, leading Dr. James Roberts and me
originally to suspect that there were closer quartzite sources; however,
we were unable to identify quartzite sources west of Draaikraal. It is
likely instead that the creeks which source their headwaters from the
Houtenbek and Steenkampsberg quartzite highlands to the east and flow
past Draaikraal to the west transported and rounded these cobbles,
which are much more stable on the surface than the RLS rocks, eventually
coming to dominate the regional alluvium.

\subparagraph{IS: Ironstone Laterites}
\label{sec:ch5-unit-is}

Ironstone laterites are common throughout the relatively low and flat
valley surrounding Draaikraal. They are somewhat weakly consolidated and
highly porous non-magnetic amalgamations of \textless1cm iron
concretions composed of 1 mm thick deep black vitreous goethite rims
surrounding pockets of fine-grained yellow limonite matrix. Interspersed
between these concretions are variably rounded medium-coarse
Steenkampsberg quartzite gravels. The goethite nodules and quartz
gravels are cemented by fine grained, red and yellow hematite-limonite
matrix with mm-scale angular, coarse, quartz sand. The existence and
abundance of these ironstone laterites implies that the valley
surrounding Draaikraal is a longstanding, persistent topographical
feature that has existed on geological timescales. They were likely
formed by in-place weathering of RLS rocks, explaining the relative
dearth of RLS material within the valley alluvium.

\subparagraph{Dl: Dolerite Sill:}
\label{sec:ch5-unit-dl}

Post-Bushveld Dolerite sills are mapped throughout the field area
\citep{vandermerwe2007,molyneux2008}, and one such discordant sill was observed
at the levels of the UMA and MMA north of the Draaikraal Lineament
within the field area (photo~\ChFivePhoto{89}). Despite being younger than the
surrounding RLS, the sill is significantly more weathered, likely
because of extensive blocky fracturing hypothesized to result from rapid
contraction during cooling. The dolerite is moderately magnetic.



\subsection{Rock Cleavage and Modal Layering}
\label{sec:ch5-cleavage-layering}

Within many of the phaneritic plutonic rocks, there was measurable,
near-planar cleavage visible (e.g. photos~\ChFivePhoto{58} and \ChFivePhoto{63}). Plutonic rock cleavages without modal layering yielded
measurements that were similar to the regional dip of bedding in the
floor metasedimentary rocks to the east, but in areas with prominent
modal layering these cleavages had strikes rotated 90 degrees from the
regional strike, interestingly dipping similarly to the exotic limb of
the Kliprivier Synform. These 90 degree offset structures could be
measured very near to one another, sometimes within the same outcrop.
Photo~\ChFivePhoto{115} shows rock cleavage that is discordant with modal layering in
a striking outcrop of the MMA.

It is important to note that one outcrop of mottled anorthosite clearly
showed disagreeing cleavage and planar modal layering, with modal layers
dipping 50° more eastwardly than the rock cleavage {[}photo~\ChFivePhoto{115}, annotated in
Figure~\ref{fig:ch5-ma-cleavage-photo}{]}.

\begin{figure}[htbp]
  \centering
  \includegraphics[width=\linewidth]{photo115_annotated_cleavage_modal_layering.jpg}
  \caption{Cleavage and modal layering forming distinct planes in a Mottled
    Anorthosite outcrop (annotated photo~\ChFivePhoto{115}).}
  \label{fig:ch5-ma-cleavage-photo}
\end{figure}

\subsection{Structural Observations}
\label{sec:ch5-structural-observations}

Structural measurements across the field area, together with the traced
lower Main Zone anorthosite marker beds, are mapped in
Figure~\ref{fig:ch5-structural-map}.

\afterpage{%
\begin{landscape}
\begin{figure}[p]
  \centering
  \includegraphics[width=\linewidth,height=0.8\textheight,keepaspectratio]{structural_map.pdf}
  \caption{Structural map of the Draaikraal field area: strike and dip of bedding,
    plutonic cleavage, and modal layering, over the mappable units (floor, Disturbed
    Critical Zone, Layered Critical Zone, Transitional Unit, Main Zone), with the lower
    Main Zone anorthosite marker beds (LMA, MMA, UMA, and the spotted anorthosites) and
    the Draaikraal Lineament (dashed). 10~m elevation contour interval. Measurements are
    listed in Table~\ref{tab:ch5-field-stations}.}
  \label{fig:ch5-structural-map}
\end{figure}
\end{landscape}}

\subsubsection{Kliprivier Synform}
\label{sec:ch5-kliprivier-synform}

\afterpage{%
\begin{landscape}
\begin{figure}[p]
  \centering
  \includegraphics[width=\linewidth,height=0.8\textheight,keepaspectratio]{kliprivier_synform_map.pdf}
  \caption{Geological map of the Kliprivier Synform. Field observations are colored by
    lithology over the mappable units (floor, Disturbed Critical Zone, Layered Critical
    Zone, Transitional Unit, Main Zone), with bedding and plutonic-cleavage measurements
    and the Draaikraal Lineament (dashed).}
  \label{fig:ch5-kliprivier-map}
\end{figure}
\end{landscape}}

The most interesting large-scale structure in the mapping area
(Figure~\ref{fig:ch5-kliprivier-map}) is a
rotated package of disturbed, pegmatoidal upper Critical Zone
pyroxenites and gabbronorites {[}photos~\ChFivePhoto{61},\ChFivePhoto{62}, \ChFivePhoto{72}--\ChFivePhoto{78}{]}, well-cleaved
leuconorite layers {[}photo~\ChFivePhoto{63}{]}, quartzites {[}photos~\ChFivePhoto{28} and \ChFivePhoto{124}{]},
migmatized/hornfelsed argillites and calc-silicate layers {[}photos~\ChFivePhoto{29},
\ChFivePhoto{30}, and \ChFivePhoto{125}{]}, and marginal sills {[}photo~\ChFivePhoto{57}{]} in the northern
portion of Kliprivier Farm. The structure is remarkable because it
apparently displaces a marginal package of metasediments and
fine-grained sills over 1 km into the plutonic RLS. This
structure is described accurately but qualitatively by
\citet{vandermerwe2007}, who investigated the area to assess its platinum resource
potential. \citet{vandermerwe2007} refers to this feature the ``Kliprivier
Syncline'', though it is not clear if he is naming the structure or it
was named elsewhere. A displaced quartzite layer within the structure
was also observed during the field mapping of the Main Zone mottled
anorthosites by \citet{latypovetal2022} and is described in their
supplementary information as a ``Large Quartzite Xenolith?''.

The structure is referred to here instead as the ``Kliprivier Synform''
to recognize that it is not yet well understood. For instance, no
previous work to date has correlated marginal magmas and country rocks
from the two limbs of the putative syncline to demonstrate that they are
the same and that the structure is indeed a fold and not another
structure, such as a laccolithically deformed roof remnant of an RLS
subchamber or a package of floor rocks entrained in a rising magma. The
map presented here offers additional details that elucidate the nature
of this structure. Below, a series of observations of the Kliprivier
Synform is presented:

\begin{enumerate}
\def\labelenumi{\arabic{enumi}.}
\item
  To the north of the structure, there is a continuous, 90° gradual
  change in strike moving from finely-banded noritic Critical Zone,
  which strikes southwest and dips northwest like the majority of the
  mapped lithologies east of Draaikraal, to a southeastern strike
  dipping southwest. This gradual change in structure is accompanied by
  lithologic changes in the Critical Zone: Near and at the deformed
  areas, grain-size increases (photo~\ChFivePhoto{61}) and stratiform modal layering
  ceases and is replaced by a heterogeneous rock with abundant
  fine-grained Marginal Zone xenoliths, anorthosite autoliths (photo~\ChFivePhoto{62})
  and irregular clumps of pyroxene surrounded by anorthosite (photos~\ChFivePhoto{72}--\ChFivePhoto{77}). I have mapped these lithologies as ``Disturbed Gabbronorite''
  (Lithologic code: Dcg).
\item
  To the south, there is little exposure and no structural data could be
  gathered to pinpoint the reversal of the structure back to the common
  northwestern dip. However, outcrops constrain the contact between the
  Marginal Zone magmas and the porphyritic gabbronorite of the upper
  Critical Zone to occur within 10 meters. The ``disturbed'' lithologies
  of the Marginal/Critical zones present on the north side of the
  synform are not present on the south side. \citet{vandermerwe2007}
  interprets faulting with downthrows to the south at this contact;
  however, \citet{vandermerwe2007} does not map any vertical displacement of
  marker beds such as the UG2 chromitite at this fault.
\item
  Stratigraphically higher plutonic rocks are undeformed, suggesting
  that they were intruded and crystallized after the deformation forming
  the Kliprivier Synform. Interestingly, the western termination of the
  synform as mapped here occurs at or slightly below the level of the
  Merensky Reef, where a significant magma recharge event corresponding
  to the emplacement of the Main Zone is thought to have occurred
  \citep{seabrook2005,cameron1982}.
\item
  Displaced plutonic rocks to the north are stratigraphically below the
  quartzites and Marginal Zone sills, meaning that the section of sills
  and sediments serve as a roof to the RLS to the north.
\item
  Within this study, there is one strike and dip measurement immediately
  to the west of the structure, that is striking SW like the majority of
  the field area but is dipping more steeply than other rocks in the
  region. \citet{latypovetal2022} also report an unusually steeply dipping
  plutonic rock along strike.
\item
  To the east, the structure is attached to the floor of the RLS.
  Unfortunately, that area was not mapped in detail and the nature of
  the rotation from the synform's exotic dip to the common dip of the
  marginal series was not observed.
\end{enumerate}

   The calculated fold axis (Figure~\ref{fig:ch5-stereonets}b) trends roughly W5°N, plunging 15°, similarly to the
Draaikraal Lineament. The calculated obtuse bisector axial plane strikes W5°N and dips 85°.

The axial plane is similar in orientation to the WNW-ESE trending
Draaikraal Lineament that bounds the topographically high and low Main
Zone rocks. 

\begin{figure}[htbp]
  \centering
  \begin{subfigure}[t]{0.48\linewidth}
    \centering
    \includegraphics[width=\linewidth]{kliprivier_stereonet_bedding_poles.jpg}
    \caption{}
    \label{fig:ch5-stereonets-bedding}
  \end{subfigure}\hfill
  \begin{subfigure}[t]{0.48\linewidth}
    \centering
    \includegraphics[width=\linewidth]{kliprivier_stereonet_fold_axis.jpg}
    \caption{}
    \label{fig:ch5-stereonets-fold-axis}
  \end{subfigure}
  \caption{Structural analysis of the Kliprivier Synform in stereographic projection
    (Stereonet 11; \citealp{allmendinger2012,cardozo2013}). (a) Bedding planes and poles
    to planes of rocks within the Kliprivier Synform. (b) Calculated fold axis and axial
    plane for the Kliprivier Synform.}
  \label{fig:ch5-stereonets}
\end{figure}

\section{Discussion}
\label{sec:ch5-discussion}

\subsection{Origin of the Kliprivier Synform and Draaikraal Lineament}
\label{sec:ch5-origin-of-the-kliprivier-monocline-and-draaikraal-lineament}

\citet{vandermerwe2007} described a ``synformal structure'' on Kliprivier
farm and referred to it as the ``Kliprivier Syncline'', but I identified
only a single, exotically-dipping limb within the structure with a
gradual southerly change in strike, therefore I interpret it here as a
monocline.

\citet{vandermerwe2007} brings attention to the PGE grade of such
structures along the eastern margin of the RLS, which often are rich in
sulfides and contain unusually high-grade Pt ores aside from the
traditional targets: the UG2 chromitite layer and the Merensky Reef 
(\S\ref{sec:ch5-the-origin-of-the-merensky-reef-and-pge-mineralization}, ``\nameref{sec:ch5-the-origin-of-the-merensky-reef-and-pge-mineralization}''). 
\citet{vandermerwe2007} believes that the Kliprivier structure and these other
sulfide traps predate Bushveld magmatism; however, I identified multiple
coarse-grained plutonic rocks of the upper Critical Zone that were rotated
within the structure. No Main Zone or Transition Zone rocks
exposed on the surface to the west of the Kliprivier are deformed within
the Kliprivier Synform, suggesting that deformation preceded the
crystallization of the Main Zone. This relative ordering requires that
the structure is synmagmatic. This interpretation is also supported by 
evidence for viscous deformation of gabbroic pegmatoids on the synform 
(Figure~\ref{fig:ch5-photo60-deformation}).

I interpret this
unit to be formed by synmagmatic deformation of nascent layered upper
Critical Zone mushes during the formation of the Kliprivier Synform,
effectively breaking up and redistributing the nascent layering into
many blobs.


% Photo 60 + digitized deformation-history sketches (2026-10-27). Schematic:
% Chapter5/Figures/fig_photo60_deformation.py -> generated/photo60_deformation_schematic.pdf
\begin{figure}[tbp]
  \centering
  \begin{subfigure}[t]{\linewidth}
    \centering
    \includegraphics[width=\linewidth]{Chapter5/Figures/generated/field_photos/DrBV60.jpg}
    \caption{}
    \label{fig:ch5-photo60-deformation-photo}
  \end{subfigure}\\[4pt]
  \begin{subfigure}[t]{\linewidth}
    \centering
    \includegraphics[width=\linewidth]{photo60_deformation_schematic.pdf}
    \caption{}
    \label{fig:ch5-photo60-deformation-schematic}
  \end{subfigure}
  \caption{(a) Sheared magnetite pegmatoid inclusion within non-magnetic marginal rocks,
    interpreted as reacted sediments, on
    the western side of the Kliprivier Synform (photo~\ChFivePhoto{60}). The
    pegmatoid is about 15~cm long. (b) Interpreted deformation history: 
    (1) mafic pegmatitic melt accumulates in lenses in
    migmatized sediment; (2) synmagmatic viscous deformation produces a sigmoid
    geometry; (3) post-magmatic brittle failure along an annealed fault truncates
    the lens and displaces its tip (lightly outlined). Black arrows show the sense
    of shear at each stage. The second stage of deformation may have occurred during isostatic flexure during cooling.
    The anneal of the fault suggests that it occurred at warm subsolidus conditions.}
  \label{fig:ch5-photo60-deformation}
\end{figure}

Structural analysis of the synform is consistent with gentle folding
along a fold axis plunging shallowly to the west with a subvertical
axial plane trending roughly E(SE)-W(NW) bisecting the limbs' obtuse
intersection angle (Figure~\ref{fig:ch5-stereonets}b). This deformation is subtle and would not
reverse the original ordering of the layered series and the marginal
series. A syncline alone cannot explain how the Marginal Zone of the
deformed, SW dipping limb came to be structurally higher than some of
the deformed plutonic rocks. It is, however, consistent with the idea
that much of the stratigraphically lower Critical Zone was intruded as
stacked sills from the north and terminated here, laccolithically deforming the overlying country rocks (\S\ref{sec:ch5-structure-and-relationship-to-country-rocks}). Within this paradigm, the
intruded Transvaal sediments ``displaced'' deeper within the RLS at the
Kliprivier Synform are actually roof rocks to the Critical Zone below
the level of the UG2. Given the intense melting and devolatilization of
the floor rocks in the area, it would not be surprising to find that the
thin strand of country rocks dividing the Critical Zone here beneath the
level of the UG2 would have been completely reacted to igneous minerals
and textures by reactive porous flow of magmas and hydrothermal fluids,
as has been argued for some Lower Zone and Marginal Zone rocks
\citep{yudovskaya2021}. Indeed, the ghosts of such transformed strands
may be expected at many locations throughout the complex and may have
separated sills of distinct composition or served as traps for
percolating melts or hydrothermal fluids. They may therefore be expected
especially at geochemical discontinuities within the stratigraphic
column, such as the Merensky Reef, Giant Mottled Anorthosite, Pyroxenite
Marker, Upper Mottled Anorthosite, etc.

\citet{sharpe1982} also reported that Critical Zone rocks but not
Main Zone rocks were deformed in the area, but attributed similar
small-scale folding on the margin to isostatic flexure. Isostatic
flexure has affected the modern expression of the Bushveld Complex
immensely, but a paleomagnetic study shows that isostatic flexure only
commenced within the Bushveld after cooling beneath the Curie
temperature of magnetite: 575°C \citep{hattingh1986}. This means that
flexure had to occur after termination of Bushveld magmatism and cannot explain synmagmatic
deformation. 

If the southwesterly dipping limb of the Kliprivier Synform is rotated
\textasciitilde20° about a level rotational axis trending
\textasciitilde45° NE to ``undo'' the local effects of postmagmatic
isostatic flexure that affected the RLS and its margin, the exotic limb
associated with the Kliprivier Synform would be dipping
\textasciitilde20° due south. This observation is consistent with the
hypothesis that an intrusion originating from the Thabazimbi-Murchison
Lineament to the north terminated in this area, producing a southerly
dipping monocline in its roof rocks.

The axial plane of the Kliprivier Synform is similar in orientation to the
Draaikraal Lineament (\S\ref{sec:ch5-kliprivier-synform}). Considering that the Kliprivier Synform also rests upon this
lineament, it is likely that the formation of the Kliprivier Synform is
related to the Draaikraal Lineament.

\subsection{Origin of Spheroidal Pyroxenite Aggregates (Tennis balls)}
\label{sec:ch5-origin-of-spheroidal-pyroxenite-aggregates-tennis-balls}

Spheroidal aggregates (SAs) of pyroxenite up to \textasciitilde5 cm are
known to occur as inclusions in the (gabbro)norites near the margin of
the RLS, and larger pyroxenite ``boulders'' within the Merensky Reef are
well-documented in the western limb of the RLS \citep{lee1979,smith2023,khoza2021}. The spherical agglomerates (also
called ``tennis balls'' due to their similar size) are richer in
pyroxenes, higher in An\# and pyroxene Mg\#, and richer in quartz and
sulfides than their (gabbro)noritic hosts \citep{lee1979,khoza2021}. Originally, dense groupings of these spherical aggregates were
thought to represent laterally continuous marker beds, but they cannot
be traced along strike of RLS layering for long distances relative to
the named reefs or mottled anorthosites-- instead they are most commonly
found \textless200m stratigraphically above the margin with the floor
rocks regardless of the stratigraphic position within the layered
series \citep{lee1979}. This author located isolated
pyroxenite tennisballs in porphyritic leuconorite in two places above
and below the Kliprivier Synform {[}photos~\ChFivePhoto{8}, \ChFivePhoto{79}{]}, on regions
previously mapped to contain tennisballs \citep{lee1979}. More
interestingly, this author located spheroidal aggregates of pyroxenite
in the fine-grained Marginal Zone {[}photos~\ChFivePhoto{2}, \ChFivePhoto{4}, \ChFivePhoto{96}{]}. Only very
recently have authors argued the marginal zone hosted spheroidal
aggregates at all. \citet{mitchell2023} used geochemical evidence
to claim the coarsely crystalline plutonic gabbronorites hosting the
Tennis Ball Marker (TBM), famously known on the Middelkraal farm, ought
to be considered as Marginal Zone and not as Main Zone as previously
mapped \citep{lee1979,molyneux2008,khoza2021}. Before
this work, no fine-grained marginal zone rocks had been reported to contain
pyroxenite tennis balls or spheroidal aggregates, although those
observed here in Draaikraal are much more lumpy and irregular than the
\emph{spheroidal} aggregates (SA) found in the Middelkraal TBM.

In photo~\ChFivePhoto{96}, the largest pyroxenite SA (\textasciitilde4 cm) resembles a
diapir: It has a long (\textasciitilde15 cm) and thin
(\textasciitilde\textless1 - 2 mm) pyroxenite tail developed behind it,
suggesting that it was migrating through a comparatively rigid,
stationary crystal-rich mush. The formation of a diapir suggests that
the SA was moving relative to the host rock and that the host rock was
more rigid than the pyroxenite. This requires explanation, as the
spheroidal aggregates contain more primitive pyroxenes and plagioclase
than the host matrix and are equilibrated at higher temperatures,
meaning that the pyroxenitic SAs would be more rigid than the host rocks
if one assumes the two are in thermal equilibrium.

There are three main hypotheses for formation of the spheroidal
pyroxenite aggregates in the RLS. The model of \citet{lee1979}
posits that pyroxenite aggregates are formed when a second, immiscible
fluid called the ``bridging'' liquid intrudes a magma body with
suspended pyroxene crystals. The bridging liquid preferentially ``wets''
the suspended pyroxene crystals, and entrains them as it coagulates with
other smaller, individual bubbles of the immiscible fluid entraining the
suspended pyroxenes of the host magma.

The second model of \citet{khoza2021} posits that a magma recharge event
introduces magma into a pyroxenite cumulus pile, breaking the pile into
pieces that are then subsequently entrained into the secondary magma.
This magma does not equilibrate with the pyroxenite; instead, it evenly
melts the rough edges of the pyroxenite aggregates, fashioning them into
spheres.

The third model posits that a pyroxenitic liquid or mush intrudes a
cooler crystal mush and is disaggregated into pyroxenite aggregates,
analogously to the formation of mafic enclaves in felsic plutons from
the disaggregation of intruded basalts \citep{mitchell2023,smith2023}. This is the model most compatible with the observation of
diapirism of spherical aggregates in the fine-grained marginal sill
phase, since it allows for the SAs to be formed by disaggregation of a
comparatively hot ultramafic magma intruded into comparatively cool (and
therefore rigid) mafic sills.

The diapiric movement of the pyroxenite spheroid in photo~\ChFivePhoto{96} serves as
good evidence against the second model of \citet{khoza2021}, as Khoza's model
requires the pyroxenite spheroids to be more rigid than the host magma,
so that they remain consolidated against the flowing host magma. If the
spheroidal aggregate was more rigid than the host magma, then it would
not have left a pyroxenitic tail behind it. Finally, if magma recharge
introduced into a pyroxenite cumulate pile forms SAs, then one would
expect them to be a ubiquitous feature throughout the Bushveld Complex
regardless of the distance from the margin, as the conventional and most
accepted model for the assembly of the Bushveld Complex evokes many
instances of magma recharge atop a preexisting cumulate pile \citep{holness2017,cawthorn2015}.

The model of \citet{lee1979} is interesting because secondary fluid
input both can explain the relative softness of the primitive SA
migrating through more evolved host rocks, and because it explains the
abundant quartz (\textgreater10\%: \citealp{lee1979}) in the
pyroxenite SA, which would not be expected for an ultramafic magma with
high Mg\# and An\#. Photos~\ChFivePhoto{97} and \ChFivePhoto{98} exhibit the migration of a distinct
secondary comagmatic silica-rich fluid through the marginal zone magmas.
However, \citet{mitchell2023} point out that multiple distinct
SA-rich horizons of the TBM exist but those that they report were over
100 meters from the alleged source of the fluids: the Dullstroom lavas
at floor of the magma chamber. Their model simply requires injection of
hot pyroxenite lavas into cooler mafic mushes. This author, however,
suggests that the secondary bridging liquids of Lee and Sharpe's model
could have percolated through homogeneous (gabbro)norites until
ultimately meeting a coarse, pyroxenitic horizon susceptible to
differential wetting.

The model of \citet{lee1979} offers a satisfactory explanation for
the SA's ties to the intrusion's margin. The secondary fluids that
circulated throughout the Bushveld chamber were predominantly sourced
from devolatilized and variably melted country rocks \citep[see also \S\ref{sec:ch5-stable-isotopes-of-the-rls}, ``\nameref{sec:ch5-stable-isotopes-of-the-rls}'']{schiffries1990,buick2000}, so Bushveld magmas near the margin would
be the most likely to encounter poorly mixed, newly-introduced
immiscible (or poorly miscible) fluids than magmas further from the
margin. The model of \citet{mitchell2023} also offers a plausible
explanation for the SAs' ties to the intrusion's margin, as the margin
would cool faster to create a larger temperature (and viscosity)
contrast between new pulses of pyroxenitic magmas and their host rocks.

\subsection{Correlation of the Mottled Anorthosites in the Lower Main Zone}
\label{sec:ch5-correlation-of-the-mottled-anorthosites-in-the-lower-main-zo}

Although there are three main named mottled anorthosite packages within
the lower Main Zone, there are also many other minor anorthosites
throughout the Main Zone \citep{sacs1980,vongruenewaldt1985}.
There is agreement between \citet{vandermerwe2007}, \citet{setera2018},
and \citet{latypovetal2022} about the location of the UMA on the
Roossenekal-Lydenburg road. I allow this agreement to calibrate the
naming of the area's mottled anorthosites. Additionally, \citet{vandermerwe2007}
and \citet{latypovetal2022} also mapped the MMA and LMA along the
Roossenekal-Lydenburg road at the same location as I have.

\citet{vandermerwe2007} maps the surface exposure of the UMA, MMA, and LMA
throughout the southern sector of the eastern limb of the RLS, and later
worked with \citet{latypovetal2022}, who undertook 6 stratigraphic
traverses throughout the southeast RLS to prove that these 3 marker beds
are laterally extensive throughout the entire Main Zone and do not abut
the country rocks, arguing that these continuous mottled anorthosite
packages were emplaced at the floor of an entirely molten 5 km thick
magma chamber. Three of these traverses overlapped with the field area
of this work. Within the field area, the UMA and MMA mapped by
\citet{vandermerwe2007} and \citet{latypovetal2022} are discordant; they do not
follow the regional dip to the northwest. I did locate anorthosites at
the coordinates provided in the SI by \citet{latypovetal2022}; however, I
conclude that the UMA and MMA that they have mapped in their traverse I,
north of the Draaikraal Lineament, are not the same anorthosites as
their mapped UMA and MMA south of the Draaikraal Lineament on the
Roossenekal-Lydenburg road in their traverse III. I further argue that
what they have mapped as UMA and MMA north of the Draaikraal Lineament
are actually spotted anorthosites and some lack mottles entirely.

Two claims should be distinguished here. The first, which I hold with
confidence, is that the anorthosites correlated with the UMA and MMA north
of the lineament by \citet{vandermerwe2007} and \citet{latypovetal2022} are
spotted anorthosites: they are thin, non-magnetic, and lack the coarse
oikocrystic mottles and associated pyroxenites that define the mottled
anorthosites of the lower Main Zone (Table~\ref{tab:ch5-ma-correlation}). The
second, which is tentative, is which of the northern mottled-anorthosite
packages corresponds to the UMA and which to the MMA.

Though this particular contested
correlation by \citet{latypovetal2022} is of minor significance to the
overall argument put forward in their work, it is essential evidence to
understand the origin and significance of the Draaikraal Lineament as
the surface expression of a synmagmatic monocline formed by laccolithic
deformation during the intrusion of much of the Critical Zone (see the end of this subsection).

The outcrops behind both correlations are compared in Table~\ref{tab:ch5-ma-correlation}; their photographs are in Appendix~\ref{app:draaikraal-field-data}.

% =============================================================
%  chapter5_ma_correlation.tex -- comparison of the lower Main Zone mottled
%  anorthosite correlations across the Draaikraal Lineament (Table). \input{} by Chapter5/Writing/chapter5_body.tex in the Discussion
%  (Correlation of the Mottled Anorthosites in the Lower Main Zone).
%
%  2026-10-29: the six-photo mosaic that followed the table was removed;
%  photos 49, 86, 87, 90 and 109 are back in the Appendix A album (with their
%  album captions), and photo 115 links to its annotated figure
%  (fig:ch5-ma-cleavage-photo).
% =============================================================

\begin{table}[tbp]
  \centering
  \footnotesize
  \setlength{\tabcolsep}{3pt}
  \caption{Lower Main Zone anorthosite outcrops compared across the Draaikraal
    Lineament. The roadcut of the UMA South of the lineament in particular is well documented and agreed to be a lower 
    Main Zone mottled anorthosite \citep{setera2018,latypovetal2022}. The
    outcrops north of the lineament as correlated by \citet{latypovetal2022} and \citet{vandermerwe2007} are compared with those in this work.
    Heights are stratigraphic heights above the Merensky Reef
    (\S\ref{sec:ch5-stratigraphic-column-methods}).}
  \label{tab:ch5-ma-correlation}
  \begin{tabular}{@{}>{\raggedright\arraybackslash}p{3.2cm}
                  >{\raggedright\arraybackslash}p{1.3cm}
                  >{\raggedright\arraybackslash}p{1.3cm}
                  >{\raggedright\arraybackslash}p{1.3cm}
                  >{\raggedright\arraybackslash}p{0.9cm}
                  >{\raggedright\arraybackslash}p{4.6cm}@{}}
    \toprule
    Outcrop & Photos & Height (m) & Thick. & Mag. & Texture and associated lithologies \\
    \midrule
    \multicolumn{6}{@{}l}{\textit{Upper Mottled Anorthosite (UMA)}}\\
    \addlinespace[2pt]
    Roadcut UMA, South
    \newline{\scriptsize $-$25.2274, 30.05494}
      & \ChFivePhoto{49}, \ChFivePhoto{53}, \ChFivePhoto{54}
      & 564 m & $\sim$0.5~m & yes
      & Anorthosite with $\sim$2cm oikocrystic pyroxene mottles overlying $\sim$0.5~m coarse feldspathic pyroxenite
        with anorthositic blocks, over fine-grained green gabbronorite with $\sim$ 1cm blebby bottle-green
        CPX; irregular, non-planar contacts \\
    \addlinespace[6pt]
    ``UMA'' of \citet{latypovetal2022}, North
    \newline{\scriptsize $-$25.1767, 30.07222}
      & \ChFivePhoto{86}
      & 1030 m & $\sim$10~cm & no
      & Anorthosite with $<$1~cm equant pyroxenes. No other unusual lithologies observed. Planar, concordant contacts \\
    \addlinespace[6pt]
    ``UMA'' of this work, ``LMA'' of \citet{latypovetal2022}, North
    \newline{\scriptsize $-$25.1821, 30.09749}
      & \ChFivePhoto{90}
      & 398 m & >10~cm & yes
      &  Anorthosite with coarse oikocrystic pyroxene mottles over coarse pyroxenite \\
    \midrule
    \multicolumn{6}{@{}l}{\textit{Middle Mottled Anorthosite (MMA)}}\\
    \addlinespace[2pt]
    Valley outcrop MMA, South
    \newline{\scriptsize $-$25.227, 30.06419}
      & \ChFivePhoto{103}--\ChFivePhoto{109}
      & 373 m & > 2 m & yes
      & Very thick; layers of oikocrystic mottled anorthosite and pyroxenite
        with both planar and wavy, discordant contacts \\
    \addlinespace[6pt]
    ``MMA'' of \citet{latypovetal2022}, North
    \newline{\scriptsize $-$25.1827, 30.08279}
      & \ChFivePhoto{87}
      & 708 m & not measured & no
      & Anorthosite with pyroxene spots and foliated (sheared) pyroxenes, no pyroxenites. \\
    \addlinespace[6pt]
    ``MMA'' of this work, North
    \newline{\scriptsize \todo{coordinates}}
      & \ChFivePhoto{115}
      & 333 m & > 4 m  & yes
      & Thick package of anorthosites with coarse oikocrystic pyroxene mottles and pyroxenites with sharp lower contacts
      that grade upward to anorthosite in $\sim$ 20 cm layers; cleavage discordant with modal layering \\
    \bottomrule
  \end{tabular}
\end{table}



The UMA outcrop on the Roossenekal-Lydenburg road (photos~\ChFivePhoto{53}, \ChFivePhoto{54})
contains a distinct package of a striking magnetic mottled anorthosite,
\textasciitilde0.5m thick underlain by a \textasciitilde0.5m
coarse-grained feldspathic pyroxenite with some reacted anorthositic
blocks and pure anorthosite bands. This pyroxenite is in turn underlain
by a unique rock interpreted to be a fine-grained gabbronorite
containing blebby anhedral bottle-green clinopyroxene phenocrysts. Fresh
faces of this rock are a uniquely vivid green color. I was able to
locate the vivid green fine grained gabbronorite with blebby anhedral
pyroxene phenocrysts (Fgog) above the Draaikraal Lineament near a cliff
with a view of the Booysendal PGE mine to the east (Fgog not pictured,
photos~\ChFivePhoto{88} and \ChFivePhoto{90}). I correlate this outcrop with the UMA south of the Draaikraal Lineament
because it preserves the same sequence of fine-grained green gabbronorite,
coarse feldspathic pyroxenite and magnetic mottled anorthosite. It sits
\textasciitilde170~m lower in the stratigraphic column (398~m versus 564~m
above the Merensky Reef; Table~\ref{tab:ch5-ma-correlation}). The match is not
perfect, but the column assumes a constant strike and dip
(\S\ref{sec:ch5-stratigraphic-column-methods}). Over the \textasciitilde6.5~km
separating the two outcrops along strike, small discrepancies between the
modeled and true strike and dip could produce significant variation in the
stratigraphic heights.

It is possible that this northern package is instead the MMA. That would
imply that the fine-grained green gabbronorite is not unique to the UMA but
also occurs beneath the MMA. It would not change the argument below: either
way, the northern correlatives are magnetic mottled anorthosites with
associated pyroxenites, and they lie at heights close to those of the
southern packages.

The anorthosite mapped by \citet{vandermerwe2007} and \citet{latypovetal2022}
as the UMA north of the Draaikraal Lineament (photo~\ChFivePhoto{86}) is very small
(\textasciitilde10 cm), non-magnetic, and lacks the associated
coarse pyroxenite and fine-green gabbronorite lithologies; it does not
resemble the UMA outcrop on the Roossenekal-Lydenburg road. Furthermore,
it is significantly stratigraphically higher: \textasciitilde1030~m above the Merensky Reef. \citet{vandermerwe2007} does not report
displacement of the UG2 over the Draaikraal Lineament, so it is
interesting that their UMA is given \textasciitilde470~m of stratigraphic displacement.

Similarly, the discordant MMA correlation they offer does not stand up.
The MMA is supposed to be a striking and very large marker in the lower
Main Zone (photos~\ChFivePhoto{103}--\ChFivePhoto{109}), and yet the outcrop they correlate (photo~\ChFivePhoto{87}) is a small exposure with subtle, non-magnetic pyroxene
spots/mottles, and no associated blebby or coarsely layered pyroxenites.
I instead find the concordant correlation of a very large package of
magnetic mottled anorthosites with coarsely layered pyroxenites at photo~\ChFivePhoto{115} much more plausible.

Neither correlation places the anorthosites at exactly the same height on
both sides of the lineament, and neither needs to. The thickness of the
anorthosite packages and of the layers between them may vary along strike,
shifting their relative positions coherently even without invoking a fault.
Nonetheless, I prefer the simplest explanation that explains the data. My
correlation follows petrography, and happens to require the smallest lateral
change to explain the observations: the mottled anorthosites can be followed
across the Draaikraal Lineament without invoking a fault.

Because there are essentially no detailed studies of the lower Main Zone
anorthosites \citep[similar to][for the GMA of the
transitional unit]{mitchell2002}, it is easy to fall prey to one's confirmation bias:
one can find an anorthosite in the lower Main Zone without too much
trouble. Detailed and careful petrography of these marker beds is
necessary to make correct correlations. The strategy of
\citet{latypovetal2022} to use these marker beds to understand the structure and
petrogenesis of the RLS is excellent; however, I would argue instead
that one should first focus on one particularly remarkable anorthosite
package (perhaps the MMA or UMA), and should follow it along strike. For
this field area, I would want to do a detailed characterization of the
Transitional unit, LMA, MMA, and UMA stratigraphy to the north of the
Draaikraal Lineament to get complete sections of the remarkable, but
incomplete, outcrops exposed along the Roossenekal-Lydenburg road. Then
I would follow the anorthosites up north to examine how their strike
might change, as their regionally concordant NW dip would quickly become
discordant to the north. I consider the following open questions to
remain:

1) Are the LMA, MMA, and UMA generally planar features, irregular
contacts aside?

2) For how long are they laterally continuous?

3) Where and why are they discordant to the regional layering, if ever?

4) Do they abut the floor?

5) How do the distances between them vary at different locations within
the intrusion?

\subsection{Origin of the Mottled Anorthosites in the Lower Main Zone}
\label{sec:ch5-origin-of-the-mottled-anorthosites-in-the-lower-main-zone}

Surprisingly little attention has been given to the spectacular
anorthositic sequences of the lower Main Zone. More often, anorthosites
from the Critical (or Transitional) Zone are studied \citep{mitchell2002,maier2016,mitchell2019,maier2020,barnes2021}, or at the top of the Main Zone at the
Pyroxenite Marker \citep{bourdeau2021,bourdeau2023}. The Upper Mottled
Anorthosite (UMA) and Pyroxenite Marker of the Main Zone are known to
demarcate significant geochemical and isotopic changes in the RLS
stratigraphy \citep{klemm1985,kruger1994,vantongeren2013}, most often ascribed to magma recharge (\S\ref{sec:ch5-radiogenic-isotope-stratigraphy-of-the-rls}, ``\nameref{sec:ch5-radiogenic-isotope-stratigraphy-of-the-rls}'').

In the East and West Limbs, the mottles of the lower Main Zone mottled
anorthosites (LMA, MMA, UMA) are surprisingly iron-rich-- the pyroxenes
have much lower Mg\# than the surrounding gabbronorites \citep{scoonmitchell1994,mitchelleales1998}. Though I was surprised to find
the mottles of these three units were magnetic at first, it had already
been known that they produce subtle but distinct magnetic anomalies
observable in aerial surveys, at least within the West Limb \citep{viljoen1985,cole2013}. \citet{viljoen1985} point out that
the stratigraphic magnetic anomalies of the mottled anorthosites of the
lower Main Zone are not present near structural weaknesses, like faults.

\citet{scoonmitchell1994} make the provocative argument that the
silicate iron-rich ultramafic pegmatites (IRUPs) of the upper Critical
and lower Main Zones are sourced from accumulated evolved fluids
interstitial to these mottled anorthosites of the lower Main Zone. These
putative residual melts are argued to flow laterally, eventually
exploiting structural weaknesses to percolate downward to replace
plagioclase-rich, magnesian cumulates with olivine-pyroxene Iron-Rich
Ultramafic Pegmatites (IRUP), a lithology where plagioclase is
conspicuously absent. This bold argument seems unnecessary. It is easy
to explain, as \citet{scoonmitchell1994} also do, that IRUPs likely
originated from downward-percolating immiscible Fe-oxide melts from the
Upper Zone. Strontium isotopes within IRUPs are uniform regardless of
stratigraphic height and are identical to the Upper Zone \citep{scoonmitchell1994}. IRUPs are known to occupy any location along the
stratigraphic column between the Upper Zone and their lowest occurrences
in the upper Critical Zone, and form a geochemical continuum that
records them becoming more magnesian and silica-rich with depth,
transitioning from Fe-Ti oxide pegmatites to Fe-rich silicate
pegmatites, telling a clear story of evolution/hybridization of
downward-percolating, dense Fe-rich liquids.

Still, the strontium isotopes of these mottled anorthosites and some
stratigraphically proximal gabbronorites are often discontinuous with
the otherwise homogeneous, isotopically evolved Main Zone immediately
surrounding them, and instead are, remarkably, identical to the Upper
Zone and IRUPs \citep{karykowski2017,kruger1994}. It may very well
be true that the IRUPs are related to the mottled anorthosites of the
lower Main Zone, though I would instead propose that these anorthosites
might have served as porous ``aquifers'' that trapped IRUP melts
originating from the Upper Zone, rather than pose the mottled
anorthosites as the source of these iron-rich melts.

Additionally, compelling geochemical arguments have been presented on
texturally similar mottled anorthosites in the Critical Zone that argue
for formation by reactive porous flow of a hydrothermal fluid including
the presence of evolved but Ni--rich olivine, hydrous phases and
sulfides, and extreme reverse zoning of plagioclase surrounding mottles \citep[see also \S\ref{sec:ch5-origin-of-enigmatic-monomineralic-igneous-rocks}, ``\nameref{sec:ch5-origin-of-enigmatic-monomineralic-igneous-rocks}'']{maier2020}. The presence of water is known to favor olivine
crystallization at the expense of pyroxenes, as well as favoring more
anorthitic plagioclase over albite \citep{baker2019}.

I envisage mottled, stratiform anorthosites to form where
low-permeability structural traps collect upward-percolating reactive
hydrothermal fluid that locally favors plagioclase crystallization and
pyroxene dissolution. EPMA work on the mottles of the LMA, MMA, and UMA
would bring clarity to the relative importance of evolved interstitial
fluids, potential IRUP contaminants, and hydrothermal fluids on the
formation of mottled anorthosites in the lower Main Zone and to
stratiform anorthosites in general.

\subsection{Rethinking the Marginal Zone: Arguments for Incremental Assembly}
\label{sec:ch5-marginal-zone-rethink}

The Marginal Zone has been traditionally understood to be the chilled
margin of the RLS \citep{harmer1985}. However, detailed work
reveals that the lowest stratigraphy of the RLS is complex. The Marginal
Zone is not always present at the floor of the RLS, especially to the
southeast \citep{molyneux2008,vandermerwe2007,latypov2020}.
Marginal Zone norites may locally intrude the Lower Zone \citep{wilson2017,yudovskaya2021}, and the Marginal Zone overlies an
additional Bushveld-related ultramafic sequence referred to as the
Apiesdooringdrai peridotite body or Basal Ultramafic Sequence (BUS) at
Burgersfort that is more likely to contain chilled components of the
earliest RLS-related intrusions \citep{wilson2012,wilson2017}. The
BUS, Marginal Zone, and Lower Zone all contain reacted sedimentary
``rafts'' or xenoliths with igneous mineralogy \citep{yudovskaya2021}.
Thinner ultramafic and mafic sills likely related to Bushveld
magmatism intrude the floor rocks kilometers below the mapped Marginal
Zone \citep{sharpehulbert1985}. In my field work, I observed that
Marginal Zone sills seemed to concordantly intrude the Transvaal
sediments below the RLS for many kilometers.

The lack of a cohesive igneous stratigraphy and the geochemical and
isotopic variability of the Marginal and Lower Zones led
\citet{mitchell2023} and \citet{scoon2023} to propose that the Marginal
and Lower Zones were assembled incrementally by magmas of different compositions (see also \S\ref{sec:ch5-magma-chamber-scales}, ``\nameref{sec:ch5-magma-chamber-scales}''). I would extend their arguments to the entire RLS. I posit
that a series of stacked peridotites, pyroxenites, norites,
gabbronorites, and gabbros were emplaced, somewhat self-sorted according
to density and the compositional control on stalling depth (more
primitive magmas will crystallize at higher temperatures) at the
location of the BUS, RLS and Marginal Zones presently. I suggest that
postcumulus modification by reactive porous flow of late stage magmas
and/or hydrothermal fluids allowed for dissolution and recrystallization
to produce phaneritic textures within the RLS, and that these fluids
interacted with the Marginal Zone to a far lesser extent than the
greater RLS. These fluids reequilibrated and reacted with the RLS within
a depth-controlled thermal gradient which produced the RLS-scale
cohesive geochemical trends that are mostly uncorrelated with isotopic
and lithological boundaries \citep[see][]{eales1986}. Subconcordant
intrusion of magmas as sills into the country rock results in the
inclusion of abundant but comparatively volumetrically small strands of
sedimentary rock in the Lower, Marginal, Critical, and Main Zones of the
RLS. These strands react to igneous mineralogy by reactive porous flow
along grain boundaries and develop planar modal layering due to lateral
transit of fluids and evolved magmas sourced from the
Thabazimbi-Murchison Lineament (TML) and they serve as traps for
hydrothermal fluids that locally supported the reaction, dissolution,
and crystallization of mafic lithologies into mottled anorthosites
below. The Merensky Reef represents a structural trap for a reduced
supercritical carbon-based fluid \citep[see also \S\ref{sec:ch5-unusual-volatile-phases-in-the-lower-rls}, ``\nameref{sec:ch5-unusual-volatile-phases-in-the-lower-rls}'']{ballhaus1985}. The
isotopic homogeneity of the Upper Zone (\S\ref{sec:ch5-open-or-closed-system}, ``\nameref{sec:ch5-open-or-closed-system}'') allows for a more orthomagmatic
interpretation: perhaps it represents a magma chamber where evolved
tholeiitic magmas and hydrothermal fluids collected, and
magnetite/anorthosite layers were crystallized by cycles of
depressurization, cooling, and perhaps magma recharge.

\subsection{Significance of Rock Cleavage and Modal Layering Orientation in the RLS}
\label{sec:ch5-cleavage-significance}

The discordance of rock cleavage and modal layering (\S\ref{sec:ch5-cleavage-layering}; Figure~\ref{fig:ch5-ma-cleavage-photo}) suggests that cleavage and modal
layering are formed by distinct processes that may have occurred at
different times, with modal layering likely predating the formation of
rock-cleavage.



There may have been multiple processes generating superimposed fabrics
in the plutonic RLS. For instance, rock fabrics may have been developed
in two phases, perhaps a primary fabric (modal layering) developed by
gravity and crystal settling (\S\ref{sec:ch5-layered-mafic-intrusions}), then a second by lateral flow of
hydrothermal fluids, or by differential contractions during cooling.
Care must be taken when interpreting planar measurements in plutonic
rocks of the RLS. The bimodal planes of the RLS in this field area may
be due to an initial development of south-dipping igneous fabric above
the Kliprivier Synform as magmas filled a local basin towards the
south before another generation of paleohorizontal fabric developed
perpendicular to thermal gradients due to differential contraction
during cooling.

\section{Conclusions}
\label{sec:ch5-conclusions}

\begin{enumerate}
\def\labelenumi{\arabic{enumi})}
\item
  The field data of this study support the interpretation that the
  Kliprivier Synform is a monocline synmagmatic with the Critical Zone
  below UG2, predating the UCZ at UG2 and up. There is no displacement
  of marker beds above UG2, including all mottled anorthosites. The
  country rocks and Marginal Zone sills of the Kliprivier Synform are
  roof rocks to the Critical Zone below UG2 and are the floor to
  undeformed Transitional and Main Zone.
\item
  The Draaikraal Lineament is best explained as the surface projection
  of the subvertical axial plane of the Kliprivier Synform. It is
  hypothesized that trapped fluids from the emplacement of the
  underlying Critical Zone may have exploited the deformed Kliprivier
  Synform, developing a vertical, planar fluid conduit at the
  Draaikraal Lineament where upwardly migrating fluids may have escaped,
  developing recessive lithologies. The regional recessive topography of
  the Main Zone south of the Draaikraal Lineament is proposed to
  originate from relatively high degrees of contamination from
  devolatilized/melted underlying sediments relative to the Main Zone
  lithologies at the same stratigraphic height to the north of the
  Draaikraal Lineament.
\item
  The discordant mottled anorthosite correlation above and below the
  Draaikraal Lineament by \citet{vandermerwe2007} and \citet{latypovetal2022}
  is locally incorrect. The field work presented here instead suggests
  that the mottled anorthosite marker beds of the lower Main Zone are
  concordant and undisplaced by the Draaikraal Lineament. More careful,
  detailed field work will be necessary to ensure correct correlation of
  the many understudied anorthosites of the lower Main Zone.
\item
  Pyroxenite aggregate tennisballs are found in fine-grained Marginal
  Zone Gabbros. This supports the interpretation that the coarse-grained
  porphyritic gabbronorite host of the TBM at Middelkraal is a member of
  the Marginal Zone after \citet{mitchell2023}. It is difficult to
  explain the development of a diapiric tail on the nascent pyroxene
  tennis balls in the fine-grained Marginal Zone with the model of
  \citet{khoza2021}, although this observation allows the models of
  \citet{mitchell2023}, \citet{smith2023}, and \citet{lee1979}.
\end{enumerate}

\section{Acknowledgment}
\label{sec:ch5-acknowledgments}

Profound thanks are due to Professor James Roberts for his charismatic
introduction to the plethora of all-consuming geological problems of
``The Beast'' (the Bushveld Complex). He is also owed thanks for warmly
hosting me as a visitor to the University of Pretoria in 2023, for
providing resources to explore the Draaikraal area, for accompanying me
for several days of field work in Draaikraal, and for teaching me
alongside UP undergraduates on a field course of the wondrous Bushveld.

